\clearpage
\fancyhead[L]{\small INHALT.}
\SourceLeaf{666}
\begingroup
\raggedright
\noindent\hangindent=1em [\textsc{Graham}'s Planet beobachtet] \mbox{.\ .\ .\ .\ .\ .\ .\ .\ .\ .\ .\ .\ .\ .\raisebox{.5ex}{\rotatebox{25}{-}}\ .\ .\ .} \hfill\mbox{1848 Juni 10}\quad\mbox{Seite 477}\par\smallskip
\noindent\hangindent=1em [\textsc{Graham}'s Planet beobachtet. Elemente von \textsc{Goldschmidt}] \mbox{.\ .\ .\ .\ .\ .} \hfill\mbox{1848 Juni 29}\quad\mbox{--- 477}\par\smallskip
\noindent\hangindent=1em [\textsc{Hind}'s Stern beobachtet] \mbox{.\ .\ .\ .\ .\ .\ .\ .\ .\ .\ .\ .\ .\ .\ .\ .\ .\ .\ .} \hfill\mbox{1848 Nov. 2}\quad\mbox{--- 478}\par\smallskip
\noindent\hangindent=1em [Parthenope beobachtet] \mbox{.\ .\ .\ .\ .\ .\ .\ .\ .\ .\ .\ .\ .\ .\ .\ .\ .\ .\ .} \hfill\mbox{1850 Juni 20}\quad\mbox{--- 478}\par\smallskip
\noindent\hangindent=1em [Victoria beobachtet] \mbox{.\ .\ .\ .\ .\ .\ .\ .\ .\ .\ .\ .\ .\ .\ .\ .\ .\ .\ .\ .} \hfill\mbox{1850 Nov. 22}\quad\mbox{--- 478}\par\smallskip
\noindent\hangindent=1em [Sonnenfinsterniss beobachtet] \mbox{.\ .\ .\ .\ .\ .\ .\ .\ .\ .\ .\ .\ .\ .\ .\ .\ .} \hfill\mbox{1851 Aug. 28}\quad\mbox{--- 479}\par\smallskip
\noindent\hangindent=1em [Psyche] \mbox{.\ .\ .\ .\ .\ .\ .\ .\ .\ .\ .\ .\ .\ .\ .\ .\ .\ .\ .\ .\ .\ .\ .\ .\ .} \hfill\mbox{1852 Juli 27}\quad\mbox{--- 479}\par\smallskip
\bigskip\noindent\textit{Beurtheilungen und Anzeigen nicht eigner Schriften.}\par\medskip
\noindent\hangindent=1em \textsc{Rohde}: Ueber die Massen der Planeten \mbox{.\ .\ .\ .\ .\ .\ .\ .\ .\ .\ .\ .\ .\ .} \hfill\mbox{1805 October}\quad\mbox{--- 483}\par\smallskip
\noindent\hangindent=1em Ueber \textsc{La Place}'s Ausdruck für Höhenmessungen durch Barometer \mbox{.\ .\ .} \hfill\mbox{1806 Mai 18}\quad\mbox{--- 485}\par\smallskip
\noindent\hangindent=1em Nachlass: von \textsc{Lindenau} an \textsc{Gauss} \mbox{.\ .\ .\ .\ .\ .\ .\ .\ .\ .\ .\ .\ .\ .\ .} \hfill\mbox{1806 Febr. 17}\quad\mbox{--- 486}\par\smallskip
\noindent\hangindent=1em Nachlass: von \textsc{Zach} an \textsc{Gauss} \mbox{.\ .\ .\ .\ .\ .\ .\ .\ .\ .\ .\ .\ .\hspace{.6em}\ .\ .\ .} \hfill\mbox{1807 April 6}\quad\mbox{--- 486}\par\smallskip
\noindent\hangindent=1em Auszug aus einem Schreiben [Ueber \textsc{Güssmann}: Cometenbahnen] \mbox{.\ .\ .\ .} \hfill\mbox{1807 März 10}\quad\mbox{--- 486}\par\smallskip
\noindent\hangindent=1em \textsc{De Zach}: Tabulae speciales aberrationis et nutationis in ascensionem rectam et in declinationem ad supputandas stellarum fixarum positiones, sive apparentes, sive veras, una cum insigniorum 494 stellarum zodiacalium catalogo novo in specula astronomica Ernestina ad initium anni 1800 constructo, cum aliis tabulis eo spectantibus \mbox{.\ .\ .\ .\ .\ .\ .\ .\ .\ .} \hfill\mbox{1808 Jan. 2}\quad\mbox{--- 489}\par\smallskip
\noindent\hangindent=1em Connoissance des tems, ou des mouvements célestes à l'usage des astronomes et des navigateurs pour l'an 1808 publiée par le bureau des longitudes. Paris 1806. Pour l'an 1809. Paris 1807 \mbox{.\ .\ .\ .\ .\ .\ .\ .\ .\ .} \hfill\mbox{1808 April 7}\quad\mbox{--- 492}\par\smallskip
\noindent\hangindent=1em \raisebox{.4ex}{\rotatebox{-45}{-}}Antoine Cagnoli: Catalogue de 501 étoiles suivi des tables relatives d'Aberration et de Nutation. Modena 1807 \mbox{.\ .\ .\ .\ .\ .\ .\ .\ .\ .\ .\ .} \hfill\mbox{1808 Mai 14}\quad\mbox{--- 496}\par\smallskip
\noindent\hangindent=1em \textsc{Pasquich}: Rechenschaft von meinen Vorschlägen zur Beförderung der Astronomie auf der königl. Universitäts-Sternwarte in Ofen \mbox{.\ .\ .\ .\ .\ .} \hfill\mbox{1808 Mai 28}\quad\mbox{--- 497}\par\smallskip
\noindent\hangindent=1em \textsc{von} \textsc{Lindenau}: Ueber den gegenwärtigen Zustand der Sternwarte Seeberg \mbox{.} \hfill\mbox{1808 Juli 2}\quad\mbox{--- 498}\par\smallskip
\noindent\hangindent=1em \textsc{Laplace}: Exposition du système du monde. Troisième édition. 1808 \mbox{.\ .\ .} \hfill\mbox{1808 Juli 23}\quad\mbox{--- 500}\par\smallskip
\noindent\hangindent=1em \textsc{Bode}: Astronomisches Jahrbuch für das Jahr 1810. Berlin 1807 \mbox{.\ .\ .\ .\ .} \hfill\mbox{1808 Sept. 26}\quad\mbox{--- 504}\par\smallskip
\noindent\hangindent=1em \textsc{Picot}: Sur les vingt-une dernières comètes et les nouvelles planètes. Bibliothèque britannique 1808 Mai. Genf \mbox{.\ .\ .\ .\ .\ .\ .\ .\ .\ .\ .\ .} \hfill\mbox{1808 Octob. 13}\quad\mbox{--- 507}\par\smallskip
\noindent\hangindent=1em \textsc{Bouvard}: Tables astronomiques publiées par le bureau des longitudes de France. Nouvelles tables de Jupiter et de Saturne calculées d'après la théorie de M. \textsc{Laplace}, et suivant la division decimale de l'angle droit. Paris 1808 \mbox{.\ .\ .\ .\ .\ .\ .\ .\ .\ .\ .\ .\ .\ .\ .\ ,\ .\ .\ .\ .} \hfill\mbox{1808 Nov. 5}\quad\mbox{--- 508}\par\smallskip
\noindent\hangindent=1em \textsc{Jabbo Oltmans}: Conspectus longitudinum et latitudinum geographicarum, per decursum annorum 1799 ad 1804 in plaga aequinoctiali ab \textsc{Alex. de Humboldt} astronomice observatarum, calculo subjecit. Paris u. Tübingen  \hfill\mbox{1808 Decbr. 3}\quad\mbox{--- 509}\par\smallskip
\noindent\hangindent=1em \textsc{Harding}: Neuer Himmelsatlas. Erste Lieferung. 1809. Hamburg bei Perthes  \hfill\mbox{1809 April 17}\quad\mbox{--- 510}\par\smallskip
\noindent\hangindent=1em Voyage \textsc{de Humboldt} et \textsc{Bonpland}. Quatrième Partie, Astronomie et Mag,netisme. Premier volume, contenant un recueil d'observations astronomiques, d'opérations trigonométriques, et de mésures barométriques, faites pendant le cours d'un voyage aux régions équinoxiales du Nouveau-Continent, depuis 1799 jusqu'en 1803. I. Lieferung 1808. II. Lieferung 1809. Paris und Tübingen \mbox{.\ .\ .\ .\ .\ .\ .\ .\ .\ .\ .\ .\ .\ .} \hfill\mbox{1809 Juni 10}\quad\mbox{--- 512}\par\smallskip
\endgroup
